\clearpage
\begin{proofbox}
  \textbf{\textcolor{CProof}{证明方法一：把根号当作一个数（最容易理解）}}
\begin{center}\begin{tikzpicture}[font=\small,>=stealth,line join=round]
\draw[gray!45,rounded corners=3pt] (0,0) rectangle (14.2,9.4);
\node at (7.1,8.8) {整个根号作新变量：沿曲线读出对应范围};

\begin{scope}[shift={(2.3,2.5)},x=1cm,y=1.7cm]
\draw[->] (0,0)--(5.7,0) node[right]{$x$};\draw[->] (0,-.2)--(0,2.8) node[above]{$t$};
\draw[teal!80!black,thick,domain=0:2.15,samples=70] plot({\x*\x+1},\x);\draw[blue,very thick,domain=1:2,samples=50] plot({\x*\x+1},\x);
\foreach \x/\y in {2/1,5/2}{\draw[dashed] (\x,0)--(\x,\y)--(0,\y);\fill[blue] (\x,\y) circle(1.7pt);\node[below] at (\x,0){$\x$};\node[left] at (0,\y){$\y$};}
\node[below] at (1,0){$1$};\fill (1,0) circle(1.5pt);
\end{scope}
\node[align=left,anchor=west] at (8.5,5.1){$t=\sqrt{x-1}\ge0$\\[10pt] $x=t^2+1$\\[12pt] 全部：$x\ge1\iff t\ge0$\\[10pt] 限制：$2\le x\le5$\\[5pt] $\phantom{限制：}\iff1\le t\le2$};
\node[text width=13.2cm,align=center] at (7.1,.65) {粗蓝线表示额外限制后的部分；换元必须同时保留它的实际范围。};
\end{tikzpicture}
\end{center}


  \textbf{要说明什么？} 设 $t=\sqrt{mx+n}$ 不是随便改名字；要确认原来的每个 $x$ 都有对应的 $t$，而且新范围内每个 $t$ 也能找到原来的 $x$。

  \begin{enumerate}[leftmargin=*,itemsep=0.5em]
    \item 根号有意义时 $mx+n\ge0$，所以 $t=\sqrt{mx+n}\ge0$。
    \item 两边平方，得到 $t^2=mx+n$。因为 $m\ne0$，可解出
      \[
        x=\frac{t^2-n}{m}.
      \]
    \item 反过来，随便取一个 $t\ge0$，按照上式算出 $x$，就有
      \[
        mx+n=t^2,\qquad \sqrt{mx+n}=\sqrt{t^2}=t.
      \]
      因此，如果原题只有根号本身的限制，$t\ge0$ 中的每个数都确实可以取到。
    \item 若原题还给出 $x$ 的范围，就把 $x=\frac{t^2-n}{m}$ 代回那些条件，进一步缩小 $t$ 的范围。
  \end{enumerate}

  \textbf{这就证明了：} 把表达式全部改成 $t$ 后，只要正确保留 $t$ 的范围，最后能得到的函数值不会改变。

  \medskip
  \end{proofbox}
\clearpage
\begin{proofbox}
\textbf{\textcolor{CProof}{证明方法二：用学过的三角函数去掉根号（常规方法）}}
\begin{center}\begin{tikzpicture}[font=\small,>=stealth,line join=round]
\draw[gray!45,rounded corners=3pt] (0,0) rectangle (14.2,9.4);
\node at (7.1,8.8) {平方差根式是上半圆，角度从正纵轴量起};

\begin{scope}[shift={(4.2,3.1)},scale=1.4]
\draw[teal!80!black,very thick] (2,0) arc(0:180:2);
\draw[->] (-2.5,0)--(2.7,0) node[right]{$x$};\draw[->] (0,-.5)--(0,2.6) node[above]{$z$};
\draw[blue,thick] (0,0)--(1.2,1.6);\draw[dashed] (1.2,0)--(1.2,1.6)--(0,1.6);\draw[orange!90!black,->] (0,.65) arc(90:53.130102:.65);
\node at (.38,.86){$\theta$};\fill (1.2,1.6) circle(1.7pt);\node[above right] at (1.2,1.6){$P$};\node[below] at (1.2,0){$a\sin\theta$};\node[left] at (0,1.6){$a\cos\theta$};\node[below left] at (-2,0){$-a$};\node[below right] at (2,0){$a$};
\end{scope}
\node[align=left,anchor=west] at (8.7,5.2){$z=\sqrt{a^2-x^2}\ge0$\\[10pt] $x=a\sin\theta$\\[8pt] $z=a\cos\theta$\\[10pt] $-\pi/2\le\theta\le\pi/2$};
\node[text width=13.2cm,align=center] at (7.1,.65) {上半圆上的点全部取到；此角度范围保证 $\cos\theta\ge0$。};
\end{tikzpicture}
\end{center}


  \textbf{情形 A：$\sqrt{a^2-x^2}$，其中 $a>0$。}

  根号有意义要求 $-a\le x\le a$。对每个这样的 $x$，都能选出一个 $\theta\in[-\frac{\pi}{2},\frac{\pi}{2}]$，使 $x=a\sin\theta$。于是
  \[
    \begin{aligned}
      \sqrt{a^2-x^2}
        &=\sqrt{a^2-a^2\sin^2\theta}\\
        &=a\sqrt{1-\sin^2\theta}\\
        &=a\sqrt{\cos^2\theta}\\
        &=a\cos\theta.
    \end{aligned}
  \]
  最后一步成立，是因为所选的角满足 $\cos\theta\ge0$。反过来，这个角度范围内每个 $\theta$ 都给出一个合法的 $x$，所以没有漏掉原题允许的值。

  \begin{center}\begin{tikzpicture}[font=\small,>=stealth,line join=round]
\draw[gray!45,rounded corners=3pt] (0,0) rectangle (14.2,9.4);
\node at (7.1,8.8) {平方和与平方差：用直角三角形核对三角代换};

\draw[blue,thick] (1.2,3.1)--(4.5,3.1)--(4.5,5.5)--cycle;\draw (4.2,3.1)--(4.2,3.4)--(4.5,3.4);
\node[below] at (2.85,3.1){$a$};\node[right] at (4.5,4.3){$x$};\node[above,rotate=36] at (2.7,4.55){$\sqrt{x^2+a^2}$};\draw (2,3.1) arc(0:36.027:0.8);\node at (2.2,3.4){$\theta$};\node at (3.3,6.9){平方和（图示 $x\ge0$）};\node at (3.3,1.9){$x=a\tan\theta$};
\draw[orange!90!black,thick] (8.1,3.1)--(11.4,3.1)--(11.4,5.5)--cycle;\draw (11.1,3.1)--(11.1,3.4)--(11.4,3.4);
\node[below] at (9.75,3.1){$a$};\node[right] at (11.4,4.3){$\sqrt{x^2-a^2}$};\node[above,rotate=36] at (9.6,4.55){$|x|$};\draw (8.9,3.1) arc(0:36.027:.8);\node at (9.1,3.4){$\theta$};\node at (10.2,6.9){平方差（$|x|\ge a$）};\node at (10.2,1.9){$|x|=a/\cos\theta$};
\node[text width=13.2cm,align=center] at (7.1,.65) {平方和允许负角对应 $x<0$；平方差另按 $x$ 正负恢复符号。};
\end{tikzpicture}
\end{center}
\textbf{情形 B：$\sqrt{x^2+a^2}$，其中 $a>0$。}

  任意实数 $x$ 都能写成 $x=a\tan\theta$，取 $-\frac{\pi}{2}<\theta<\frac{\pi}{2}$。因为这一段 $\cos\theta>0$，
  \[
    \begin{aligned}
      \sqrt{x^2+a^2}
        &=a\sqrt{1+\tan^2\theta}
         =a\sqrt{\frac{1}{\cos^2\theta}}
         =\frac{a}{\cos\theta}.
    \end{aligned}
  \]
  这里只用了 $1+\tan^2\theta=\frac1{\cos^2\theta}$，不需要额外学习新三角函数名称。

  \textbf{情形 C：$\sqrt{x^2-a^2}$，其中 $a>0$。}

  先考虑 $x\ge a$，令
  \[
    x=\frac{a}{\cos\theta},\qquad 0\le\theta<\frac{\pi}{2}.
  \]
  由于 $\cos\theta>0$、$\tan\theta\ge0$，
  \[
    \begin{aligned}
      \sqrt{x^2-a^2}
        &=a\sqrt{\frac{1}{\cos^2\theta}-1}
         =a\sqrt{\tan^2\theta}
         =a\tan\theta.
    \end{aligned}
  \]
  如果 $x\le-a$，则改用 $x=-\frac{a}{\cos\theta}$，仍取同样的角度范围，根号也等于 $a\tan\theta$。\textbf{两边都要考虑，不能只算 $x\ge a$。}

  \medskip
  \end{proofbox}
\clearpage
\begin{proofbox}
\textbf{\textcolor{CProof}{证明方法三：只用平方差公式，不必使用三角函数（快捷方法）}}
\begin{center}\begin{tikzpicture}[font=\small,>=stealth,line join=round]
\draw[gray!45,rounded corners=3pt] (0,0) rectangle (14.2,9.4);
\node at (7.1,8.8) {共轭和始终为正，并且能取遍所有正数};

\begin{scope}[shift={(4.4,2.3)},x=.95cm,y=.9cm]
\draw[->] (-3.4,0)--(3,0) node[right]{$x$};\draw[->] (0,-.3)--(0,5.8) node[above]{$t$};
\draw[blue,thick,domain=-3.2:2.6,samples=100] plot(\x,{sqrt(\x*\x+1)+\x});\fill[blue] (0,1) circle(1.7pt);\node[above left] at (0,1){$(0,1)$};\draw[dashed,orange!90!black] (-3.2,3)--(1.3333333,3)--(1.3333333,0);\fill[orange!90!black] (1.3333333,3) circle(1.7pt);\node[below right=9pt] at (1.3333333,3){$(4/3,3)$};
\end{scope}
\node[align=left,anchor=west] at (8.3,5.1){图示 $a=1$：\\[7pt] $t=\sqrt{x^2+1}+x>0$\\[10pt] $x=(t-1/t)/2$\\[10pt] $\sqrt{x^2+1}-x=1/t$\\[10pt] 每个 $t>0$ 都有原像};
\node[text width=13.2cm,align=center] at (7.1,.65) {一般地，两共轭式之积为 $a^2$；反解 $x=(t^2-a^2)/(2t)$。};
\end{tikzpicture}
\end{center}


  \textbf{情形 A：$\sqrt{x^2+a^2}$，其中 $a>0$。}

  设 $t=\sqrt{x^2+a^2}+x$。因为 $\sqrt{x^2+a^2}>|x|$，所以 $t>0$。又因为
  \[
    (\sqrt{x^2+a^2}+x)(\sqrt{x^2+a^2}-x)=a^2,
  \]
  可得
  \[
    \sqrt{x^2+a^2}-x=\frac{a^2}{t}.
  \]
  把“加”的等式和“减”的等式分别相加、相减：
  \[
    \boxed{\sqrt{x^2+a^2}=\frac{t^2+a^2}{2t}},\qquad
    \boxed{x=\frac{t^2-a^2}{2t}}.
  \]
  反过来，对任意 $t>0$，用第二个式子算出的 $x$，都会使第一个式子成立；因此 $t$ 确实能够取遍所有正数。

  \begin{center}\begin{tikzpicture}[font=\small,>=stealth,line join=round]
\draw[gray!45,rounded corners=3pt] (0,0) rectangle (14.2,9.4);
\node at (7.1,8.8) {平方差根式有左右两支，根号值都保持非负};

\begin{scope}[shift={(4.1,2.6)},scale=1.4]
\draw[->] (-2.5,0)--(2.7,0) node[right]{$x$};\draw[->] (0,-.4)--(0,2.9) node[above]{$z$};
\draw[blue,thick,domain=0:2.1,samples=70] plot({sqrt(1+\x*\x)},\x);\draw[orange!90!black,thick,domain=0:2.1,samples=70] plot({-sqrt(1+\x*\x)},\x);
\fill[blue] (1,0) circle(1.6pt);\fill[orange!90!black] (-1,0) circle(1.6pt);\node[below] at (1,0){$a$};\node[below] at (-1,0){$-a$};
\end{scope}
\node[align=left,anchor=west] at (8.4,5.1){$z=\sqrt{x^2-a^2}$\\[10pt] 右支：$t=x+z\ge a$\\[10pt] 左支：$t=-x+z\ge a$\\[10pt] 两支均有\\[5pt] $z=(t^2-a^2)/(2t)\ge0$};
\node[text width=13.2cm,align=center] at (7.1,.65) {图形按 $a=1$ 绘制；仅给出根式定义域时，两支都不能漏。};
\end{tikzpicture}
\end{center}
\textbf{情形 B：$\sqrt{x^2-a^2}$，其中 $a>0$、$x\ge a$。}

  设 $t=x+\sqrt{x^2-a^2}\ge a$。平方差公式给出
  \[
    (x+\sqrt{x^2-a^2})(x-\sqrt{x^2-a^2})=a^2,
    \qquad x-\sqrt{x^2-a^2}=\frac{a^2}{t}.
  \]
  把两式相加、相减，得到
  \[
    \boxed{x=\frac{t^2+a^2}{2t}},\qquad
    \boxed{\sqrt{x^2-a^2}=\frac{t^2-a^2}{2t}}\quad(t\ge a).
  \]
  对任意 $t\ge a$，上述第二式非负，第一式满足 $x\ge a$，所以没有引入原题不能取的数。

  若 $x\le-a$，把上面推导中的 $x$ 换成 $-x\ge a$ 即可：令 $t=-x+\sqrt{x^2-a^2}\ge a$，得到
  \[
    x=-\frac{t^2+a^2}{2t},\qquad
    \sqrt{x^2-a^2}=\frac{t^2-a^2}{2t}.
  \]

  \medskip
  \textbf{\textcolor{CProof}{证明结束，回到解题时的三件事：}}
  \textbf{写清原题限制、换对新字母、验证新范围内的数都能取到。} 这三件事做好，换元前后所求的值域才真正相同。
\end{proofbox}
